%LTeX: language=es-AR

\documentclass[12pt]{amsart}

%Fuente
\usepackage[utf8]{inputenc}
\usepackage[spanish]{babel}
\usepackage[T1]{fontenc}
\usepackage{lmodern}
\usepackage{epigraph}
\usepackage{lipsum}
\usepackage[a4paper, margin=2.5cm]{geometry}
\usepackage{multicol}
\usepackage{amsmath}
\usepackage{amssymb}
\usepackage[shortlabels]{enumitem}
\usepackage{listings}
\lstdefinestyle{micode}{
    basicstyle=\ttfamily\small,
    breaklines=true,
    columns=fullflexible,
    keepspaces=true,
    showstringspaces=false
}

\DeclareMathOperator{\mcd}{mcd}

\renewcommand{\refname}{Bibliografía}

%diagramas
\usepackage{tikz-cd}

%enumerar páginas y encabezados
\pagestyle{headings}

%Comandos para escribir C más rápido
\usepackage{mathrsfs}
\newcommand{\CC}{\mathbb{C}}
\newcommand{\NN}{\mathbb{N}}
\newcommand{\QQ}{\mathbb{Q}}
\newcommand{\RR}{\mathbb{R}}
\newcommand{\ZZ}{\mathbb{Z}}
\newcommand{\FF}{\mathbb{F}}
\newcommand{\cont}{\mathcal{C}}

\newcommand{\pp}{\mathfrak{p}}

\newcommand{\aaa}{\mathfrak{a}}
\newcommand{\bbb}{\mathfrak{b}}

\newcommand{\OO}{\mathcal{O}}

\newcommand{\leg}[2]{\left( \frac{#1}{#2} \right)}
\newcommand{\minimat}[4]{\left(\begin{smallmatrix} #1 & #2 \\ #3 & #4 
\end{smallmatrix}\right)}
\newcommand{\lp}{\left(}
\newcommand{\rp}{\right)}
\newcommand{\lb}{\left|}
\newcommand{\rb}{\right|}
\newcommand{\lc}{\left<}
\newcommand{\rc}{\right>}
\newcommand{\lco}{\left[}
\newcommand{\rco}{\right]}

%dobles implicaciones
\usepackage{csquotes}

%Entornos con estilo 'tipo teorema'
\theoremstyle{plain}
\newtheorem{teo}[equation]{Teorema}
\newtheorem{ej}[equation]{Ejercicio}
\newtheorem{prop}[equation]{Proposición}
\newtheorem{lem}[equation]{Lema}
\newtheorem{coro}[equation]{Corolario}
\newtheorem{defi}[equation]{Definición}
\newtheorem{obs}[equation]{Observación}
\newtheorem*{nota*}{Notación}
\renewcommand*{\proofname}{Demostración}

%Hipervínculos cliqueables, en color azul
\usepackage[colorlinks=true,allcolors=blue]{hyperref}

\title{Título}

\author{Pedro Nicolás Peña}

\begin{document}

\maketitle

% con Alt + Z te acomoda las lineas

\section*{Sección}

Vamos a trabajar mucho con la curva $E_t: y^2 = x(x+1)(x+t^2)$ porque tiene grupo de torsión isomorfo a $\ZZ/2\ZZ \oplus \ZZ/4\ZZ$ sobre $\QQ$ y si $t^2-1$ es un cuadrado entonces sobre $\QQ(i)$ es $\ZZ/4\ZZ \oplus \ZZ/4\ZZ$. Otra presentación de la curva es $E_{mn}: y^2 = x(x+n^2)(x+m^2)$ con la condición siendo que $m^2-n^2$ sea un cuadrado. Y más en general la curva va a ser $E_{rs}: y^2 = x(x+r^2)(x+s^2)$ donde $r$ y $s$ van a ser polinomios en las variables $t$ o $m$ y $n$. El cambio de variables es fácil de ver: si $t = \frac{r}{s}$ entonces

\[
y^2 = x(x+1)\lp x+\leg{r}{s}^2\rp \iff {(y's^{-3})}^2 = (x's^{-2})((x's^{-2})+1)\lp (x's^{-2})+\leg{r}{s}^2\rp 
\]

\[
\iff y'^2 s^{-6} = x'(x'+s^2)(x'+r^2) s^{-6} \iff y'^2 = x'(x'+r^2)(x'+s^2)
\]


Notemos también que es la curva de Legendre $E_\lambda : y^2 = x(x+1)(x+\lambda)$ con $\lambda = t^2$, sobre la cual hay mucha literatura.

$\textbf{Los puntos:}$

Sí ponemos $P = (rs,rs(r+s))$ y $Q = (r(u-r), iru(u-r))$ con $u^2=r^2-s^2$

\[
\begin{array}{|c|c|c|c|c|}
\hline
+   & \OO    & P      & 2*P    & 3*P \\ \hline

\OO & \OO     & (rs,rs(r+s)) & (0,0)     & (rs,-rs(r+s)) \\ \hline

Q& (-r(r-u),  & (-s(s-iu), & (-r(r+u), & (-s(s+iu), \\

&  -iru(r-u)) & su(s-iu))  & iru(r+u)) & -su(s+iu)) \\ \hline

2*Q& (-r^2,0) & (-rs,-rs(r-s))& (-s^2,0)& (-rs,rs(r-s)) \\ \hline

3*Q & (-r(r-u),& (-s(s+iu, & (-r(r+u), & (-s(s-iu), \\ 

&   iru(r-u)) & su(s+iu)) & -iru(r+u)) & -su(s-iu)) \\ \hline
\end{array}
\]

(esta tabla está chequeada a mano con Sage)

$\textbf{Discriminante:}$

\[
\Delta_{Ers} = 16 r^4 s^4 {(r^2-s^2)}^2 = 2^4 r^4 s^4 u^4
\]





$\textbf{Mas puntos:}$

Sí queremos más puntos necesitamos que $x(x+r^2)(x+s^2)$ sea un cuadrado en el cuerpo (recordar que estamos en $\QQ(i)$), por ejemplo podemos pedir $\pm x = (x+r^2)(x+s^2)$ es decir $0 = x^2+(r^2+s^2\pm 1)x +r^2s^2$ es decir que ${(r^2+s^2\pm 1)}^2-4r^2s^2$ sea un cuadrado. Sí ponemos el resto de fórmulas queremos que los siguientes sean cuadrados en el cuerpo:
\[
{(r^2+s^2\pm 1)}^2-4r^2s^2 
\]
\[
{(r^2\pm 1)}^2\mp 4s^2
\]
\[
{(s^2\pm 1)}^2\mp 4r^2
\]

Obvio habría que ver que estos puntos no terminen siendo los de torsión. 

Tal vez el teorema que me dé puntos independientes no sea estilo Green-Tao sino con ternas pitagóricas...

Ejemplo:

\begin{lstlisting}[style=micode]

Qi.<i> = QuadraticField(-1)
Ee = EllipticCurve(Qi, [0, 7^2 + 24^2, 0, 7^2*24^2, 0])

P = Ee(-18, i*558)
Q = Ee(-32, i*544) 

print(2*P == -2*Q)
print(P == Ee(-24^2,0) - Q)

print(P)
print(P*2)
print(P*3)
print(P*4)
print(P*5)
print(P*6)

True
True
(-18 : 558*i : 1)
(-625 : -4200*i : 1)
(-16006482/368449 : -80192491794/223648543*i : 1)
(-131334484801/70560000 : 39027433524414049/592704000000*i : 1)
(-8670396132425682/45913408358384449 : 716734971665869348872396306/9838056228201864674002207*i : 1)
(-10263711005525248100880625/7609855338495900009601 : -24432799810728092317281273816037452600/663841809863829299588354188574401*i : 1)
\end{lstlisting}



\begin{prop}
Un punto $(x_1,y_1)$ en $E(k): y^2 = (x-\alpha_1)(x-\alpha_2)(x-\alpha_3)$ se puede dividir por $2$ si y solo si los números $x_1-\alpha_i$ son cuadrados en $k$.
\end{prop}













\end{document}